\documentclass[12pt,a4paper]{article}

\input{glyphtounicode}
\pdfgentounicode=1

\usepackage[english]{babel}
\usepackage[utf8]{inputenc}
\usepackage[T1]{fontenc}
\usepackage{lmodern}
\usepackage{amsmath,amssymb,amsthm,mathtools}
\usepackage{xurl}
\usepackage[colorlinks=true,linkcolor=blue,citecolor=red,urlcolor=blue,unicode=true]{hyperref}
\usepackage{geometry}
\geometry{margin=2.5cm}
\usepackage{titling}
\usepackage{enumitem}
\usepackage{array}
\usepackage{booktabs}
\usepackage{tabularx}
\usepackage[capitalise,noabbrev]{cleveref}

\newtheorem{theorem}{Theorem}[section]
\newtheorem{lemma}[theorem]{Lemma}
\newtheorem{proposition}[theorem]{Proposition}
\newtheorem{corollary}[theorem]{Corollary}
\newtheorem{conjecture}[theorem]{Conjecture}
\theoremstyle{definition}
\newtheorem{definition}[theorem]{Definition}
\theoremstyle{remark}
\newtheorem{remark}[theorem]{Remark}

\newcommand{\HH}{\mathbb{H}}
\newcommand{\RR}{\mathbb{R}}
\newcommand{\CC}{\mathbb{C}}
\newcommand{\NN}{\mathbb{N}}
\newcommand{\ZZ}{\mathbb{Z}}
\newcommand{\QW}{\mathrm{QW}}
\newcommand{\Sel}{Z_{\Gamma}}
\newcommand{\Delop}{\Delta_{\Gamma}}
\newcolumntype{L}[1]{>{\raggedright\arraybackslash}p{#1}}
\newcolumntype{Y}{>{\raggedright\arraybackslash}X}
\setlength{\droptitle}{-2.4em}
\pretitle{\begin{center}\Large}
\posttitle{\par\end{center}\vspace{-0.8em}}
\preauthor{\begin{center}\normalsize}
\postauthor{\par\end{center}\vspace{-0.7em}}
\predate{\begin{center}\normalsize}
\postdate{\par\end{center}\vspace{-1.2em}}

\title{NE-B Failure as Hilbert--P\'olya Detection:\\
Universality of the v2.0 Weil Quadratic Form on the Selberg Zeta Function}
\author{Lukas Geiger\\
\small Bernau, Germany\\
\small ORCID: \href{https://orcid.org/0009-0005-7296-1534}{0009-0005-7296-1534}}
\date{\today}

\hypersetup{
  pdftitle={NE-B Failure as Hilbert--P\'olya Detection: Universality of the v2.0 Weil Quadratic Form on the Selberg Zeta Function},
  pdfauthor={Lukas Geiger},
  pdfsubject={Selberg zeta function transfer paper for the v2.0 Weil quadratic form framework},
  pdfkeywords={Selberg zeta function, Weil quadratic form, Shift Parity Lemma, Hilbert--P\'olya, hyperbolic surface, Laplace operator, NE-B}
}
\makeatletter
\renewcommand*{\HyperDestNameFilter}[1]{en.#1}
\makeatother

\begin{document}
\overfullrule=0pt
\setlength{\emergencystretch}{2.5em}
\maketitle
\vspace{-1.4em}

\begin{abstract}
The v2.0 methodology recently proposed for the Riemann Hypothesis (Geiger 2026,
Part~II) reduces RH to \emph{even dominance} of the Weil quadratic form
$\QW_\lambda$ via three algebraic and analytic ingredients: the Shift Parity
Lemma, frontier-prime dominance, and two non-existence theorems (NE-A:
non-positivity of the prime shift multiplier $M(\xi) = -2\,\mathrm{Re}[\zeta'/\zeta(\tfrac{1}{2}+i\xi)]$,
NE-B: no universal commuting operator). A natural question is whether this
framework is specific to the Riemann zeta function or whether it is a
general-purpose method for explicit-formula zeta structures. We test it on the
Selberg zeta function $\Sel(s)$ of a compact hyperbolic surface
$\Gamma\backslash\HH$. We show:

\begin{enumerate}[label=(\roman*),leftmargin=*,itemsep=0.15em,topsep=0.25em,parsep=0pt,partopsep=0pt]
\item The Shift Parity Lemma transfers unchanged (it is purely algebraic).
\item Frontier dominance transfers with modified constants, accounting for the
exponential density of primitive closed geodesics $\pi_\Gamma(T) \sim e^T/T$
in place of $\pi(x) \sim x/\log x$.
\item A Selberg--Weil quadratic-form analogue $\QW^\Gamma_\lambda$ is
obtained on the trace-formula test-function side; its relation to geometric
flow operators is mediated by the Selberg transform, not by an explicit
unitary intertwiner constructed here.
\item \emph{Even dominance of $\QW^\Gamma_\lambda$}, together with the
no-exceptional-eigenvalue condition $\lambda_1(\Delop)\ge 1/4$ and an
explicit C2/minimizer-identification gate, gives a conditional v2.0
consistency reduction of the strict statement for the non-trivial spectral
zeros in the critical strip.
\item \emph{The Selberg trace-formula channel has a Casimir commutant.} In
the spherical channel, the Laplacian $\Delop$ commutes with the relevant
radial convolution algebra. Thus the Selberg case exhibits the structural
phenomenon that NE-B is meant to detect, but this paper does not identify the
truncated $L^2([-L,L])$ shift model with the full geodesic-flow representation.
\end{enumerate}

The failure of NE-B in the Selberg setting is not a defect; it reflects the
geometric coordination of geodesics on a common surface, which is absent for
the primes. We deduce a cautious meta-principle: the v2.0 framework supplies a
common quadratic-form test language, while Hilbert--P\'olya is a \emph{special
case} whose existence depends on whether the underlying counting objects admit
a hidden common symmetry. In the Selberg spherical trace-formula channel they
do; in the Riemann case no comparable commutant is supplied by the tested NE-B
route. The Weil quadratic form remains the common object, while an operatorial
realization is an additional structure.
\end{abstract}
\vspace{-0.8em}

\noindent\textbf{MSC 2020:} 11M36, 11F72, 58J50, 47A75\\
\textbf{Keywords:} Selberg zeta function, Weil quadratic form, Shift Parity
Lemma, Hilbert--P\'olya, hyperbolic surface, Laplace operator, NE-B.

\clearpage
\tableofcontents
\clearpage

% =============================================================================
\section{Introduction}
\label{sec:intro}
% =============================================================================

\subsection{Motivation and main claim}

The v2.0 programme for the Riemann Hypothesis \cite{Geiger2026partII}
reduces RH to a quadratic-form inequality. Specifically, Connes' 2026 survey
\cite{Connes2026} writes RH as the statement that the minimum eigenvalue of
a one-parameter family of quadratic forms $\QW_\lambda$ on
$L^2([-\log\lambda,\log\lambda])$ is simple and has an even eigenfunction.
The v2.0 approach establishes this in three layers:

\begin{itemize}
\item a purely algebraic \emph{Shift Parity Lemma} (Theorem~4.1 of \cite{Geiger2026partII}),
showing that the even--odd difference matrix $D_N(r)$ has strictly negative
minimum eigenvalue for all $r \in (0,2)$, $N \geq 3$;

\item an analytic \emph{frontier-dominance} mechanism: the gap
$\Delta(\lambda) = \lambda_1^+ - \lambda_1^-$ is driven by primes with
$\log p / \log\lambda$ close to $1$ (``frontier primes''), whose density is
controlled by the Prime Number Theorem;

\item two structural non-existence results that close off alternative routes:
\emph{NE-A} (the prime shift multiplier $M(\xi)$ is not nonnegative a.e.,
ruling out $\mathrm{PF}_\infty$ total positivity) and \emph{NE-B} (no nontrivial
symmetric matrix commutes with all $D_N(r)$ simultaneously, ruling out
Sturm--Liouville-type universal operators at the tested truncations).
\end{itemize}

The v2.0 programme succeeds \emph{without} constructing a Hilbert--P\'olya
operator. It is natural to ask whether this methodological package transfers
to other zeta-like functions, and, if so, what it teaches us. The paradigmatic
case is the Selberg zeta function $\Sel(s)$ of a compact hyperbolic surface
$M = \Gamma\backslash\HH$, where:

\begin{itemize}
\item the geometric explicit formula (Selberg trace formula \cite{Selberg1956})
plays the role of the Weil explicit formula \cite{Weil1952};
\item the lengths of primitive closed geodesics play the role of primes;
\item the spectral zero description is a theorem via the self-adjointness of
the Laplacian $\Delop$; the strict critical-line form additionally requires
the no-exceptional-eigenvalue condition $\lambda_1(\Delop)\ge 1/4$.
\end{itemize}

Thus we ask two concrete questions:

\medskip
\noindent\textbf{Q1.} \emph{Do the v2.0 ingredients (Shift Parity Lemma,
frontier dominance, Weil quadratic form, NE-A, NE-B) all transfer to the
Selberg setting?}

\medskip
\noindent\textbf{Q2.} \emph{If any of them fails to transfer, what is the
structural reason?}

\medskip
Our answers, in brief:

\begin{itemize}
\item Shift Parity Lemma: \textbf{transfers unchanged} (algebraic).
\item Frontier dominance: \textbf{transfers with modified constants},
reflecting $\pi_\Gamma(T) \sim e^T/T$.
\item Weil quadratic form: \textbf{natural analogue $\QW^\Gamma_\lambda$} on
the trace-formula test-function side; the full geometric-flow intertwiner is
not constructed here.
\item NE-A: \textbf{transfers}, after locating the sign changes in the
trivial/topological/identity terms rather than in the critical-line zeros.
\item NE-B: \textbf{the obstruction is absent in the spherical trace-formula
channel}. The Casimir/Laplacian $\Delop$ commutes with radial
$K$-bi-invariant convolution operators. This is the operatorial Selberg
channel; it is not an identification of individual geodesic-flow pullbacks with
the truncated interval shifts.
\end{itemize}

The absence of the relevant NE-B obstruction is the scientifically interesting
outcome. It is not a failure of v2.0; it is a \emph{diagnostic}. Where such a
commutant exists in the right function channel, the Hilbert--P\'olya programme
has an operatorial route. Where the tested NE-B route admits no non-trivial
commutant (as in the Riemann truncation model), this universal commutant
strategy is blocked. In either case the v2.0 framework remains applicable as a
quadratic-form test.

\begin{remark}[Position within the meta-framework]
This paper operates as the \emph{Lie-origin companion} to the arithmetic
Riemann case \cite{Geiger2026partII} and sits within the three-component
decomposition $\mathrm{RH}_X = \mathrm{GBC}_X + C2_X + \mathrm{Hurwitz}_X$
developed in the companion meta-paper \cite{GeigerMeta2026}. Selberg
is the paradigmatic case in which the spherical/Casimir channel supplies the
operatorial side of $C2_X$; Riemann is the case in which
the parity substep is closed by v2.1 but the eigenvector-realisation
substep remains open. See \Cref{rem:meta-paper-refinement} below for
details.
\end{remark}

\subsection{Main theorem (informal)}

\begin{theorem}[v2.0 Selberg transfer, informal]
\label{thm:main-informal}
Let $M = \Gamma\backslash\HH$ be a compact hyperbolic surface and
$\Sel(s)$ its Selberg zeta function. The v2.0 methodology applies to the
Selberg quadratic form $\QW^\Gamma_\lambda$, subject to an explicit
C2/minimizer-identification gate, and conditionally reduces the strict
statement for the non-trivial spectral zeros in the critical strip to even dominance plus
$\lambda_1(\Delop)\ge 1/4$. The analogue of NE-A holds,
but the NE-B obstruction is absent in the spherical Selberg channel: the
Casimir/Laplace operator $\Delop$ commutes with the radial convolution
operators appearing in the trace-formula representation. The resulting
route is therefore a conditional consistency test, and this exposes a
meta-principle: the Weil quadratic form is the common test object, whereas a
Hilbert--P\'olya operator is an additional structure detected by failure of
the relevant NE-B obstruction.
\end{theorem}

\subsection{Structure of the paper}

Section~\ref{sec:v2-recollect} recalls the v2.0 ingredients.
Section~\ref{sec:selberg-setting} sets up the Selberg trace formula and the
geodesic-length analogue of the explicit formula.
Section~\ref{sec:v2-to-selberg} carries out the transfer of the Shift Parity
Lemma, frontier dominance, and constructs $\QW^\Gamma_\lambda$; a conditional
critical-line reduction is derived.
Section~\ref{sec:ne-b-fails} proves the absence of the NE-B obstruction in
the spherical Selberg channel and identifies the Casimir commutant.
Section~\ref{sec:meta} draws the meta-principle: Hilbert--P\'olya is a
special case of v2.0.
Section~\ref{sec:applications} discusses applications and outlook
(generalised L-functions, Ruelle zeta, quantum chaos).
Section~\ref{sec:conclusion} concludes.


% =============================================================================
\section{Recollection of the v2.0 Method}
\label{sec:v2-recollect}
% =============================================================================

We briefly summarize the four key ingredients of the v2.0 approach to RH. A
full treatment is in \cite{Geiger2026partII}.

\subsection{The Weil quadratic form}

Set $L = \log\lambda$. On $L^2([-L,L])$, define
\begin{equation}
\label{eq:QW-riemann}
\begin{split}
\langle A_\lambda f\mid f\rangle
&= (\log 4\pi + \gamma)\|f\|^2 \\
&\quad + \int_0^\infty \Big[\langle T_u f + T_{-u} f, f\rangle - 2e^{-u/2}\|f\|^2\Big]
\frac{e^{u/2}}{2\sinh u}\,du \\
&\quad + \sum_p\sum_{m\ge 1}\frac{\log p}{p^{m/2}}
\langle T_{m\log p} f + T_{-m\log p} f, f\rangle,
\end{split}
\end{equation}
where $T_s$ is the shift operator $(T_s f)(t) = f(t-s)\mathbf{1}_{[-L,L]}(t-s)$
and $\gamma$ is the Euler--Mascheroni constant. The quadratic form $\QW_\lambda$
is the quadratic form of $A_\lambda$.  The key algebraic tool is
Theorem~6.1 of \cite{Connes2026} (proved jointly with van Suijlekom,
building on \cite{ConnesMoscovici2022}):

\begin{theorem}[{\cite[Theorem~6.1]{Connes2026}}]
\label{thm:connes}
Let $L > 0$, let $\mathcal{D}$ be a real distribution on the interval
$[0,L]$, and let $\tilde{\mathcal{D}}$ be the associated even distribution
on $[-L,L]$.  Assume that the quadratic form with Schwartz kernel
$\tilde{\mathcal{D}}(x-y)$ defines a lower-bounded selfadjoint operator on
$L^2\!\bigl([-\tfrac{L}{2},\tfrac{L}{2}]\bigr)$, and that the minimum of its
spectrum is a simple, isolated eigenvalue with even eigenfunction~$\eta$.
Then all zeros of the entire function $\tilde\eta(z)$, $z\in\mathbb{C}$
(Fourier transform of\/~$\eta$), lie on the real line.
\end{theorem}

Applied to the Weil quadratic form $\QW_\lambda$ (with $\mathcal{D}$ the Weil
distribution from the prime explicit formula), if the minimum eigenvalue is
simple with even eigenfunction for all $\lambda\ge\lambda_0$ and the
corresponding eigenvectors converge to Riemann's $\Xi$-function, then the
zeros of $\tilde\eta$ coincide with those of $\Xi$; hence RH holds.  The
convergence step remains open \cite[§6.6]{Connes2026}.

\begin{remark}[Interval convention]
\label{rem:interval-convention}
Theorem~\ref{thm:connes} places the quadratic-form operator on
$L^2\!\bigl([-\tfrac{L}{2},\tfrac{L}{2}]\bigr)$ in Connes' notation.
In the present paper (Section~\ref{sec:v2-to-selberg} onwards),
test functions live on $L^2([-L,L])$.  Applying
Theorem~\ref{thm:connes} requires the identification $L_{\mathrm{Thm}} = 2L$,
i.e.\ the half-length in Connes' statement equals the full truncation length~$L$.
\end{remark}

\begin{remark}[Independence from Connes' convergence step]
\label{rem:connes-hedge}
Theorem~\ref{thm:connes} is drawn from Connes' 2026 preprint \cite{Connes2026};
the convergence of the construction to the Weil distribution is noted as still open
\cite[§6.6]{Connes2026}.  In the Selberg application
(Section~\ref{sec:v2-to-selberg}), this open step is \emph{immaterial}: the
spectral parametrisation of the zeros is independently secured by Selberg's
classical theory via the Laplacian $\Delop$ \cite{Selberg1956}; the strict
critical-line form additionally requires the no-exceptional-eigenvalue
condition $\lambda_1(\Delop)\ge 1/4$.
Theorem~\ref{thm:connes} plays a structural consistency role, not a logically
necessary one.  The broader programme of encoding $L$-functions via spectral
quadratic forms is additionally supported by Connes' 1999 noncommutative
geometry approach \cite{Connes1999}, Deninger's regularized-determinant
programme \cite{Deninger1994}, and the Bost--Connes thermodynamical system
\cite{BostConnes1995}.
\end{remark}

\subsection{The Shift Parity Lemma}

Let $\varphi_n(t) = \cos(n\pi t/L)/\sqrt{L}$ (even basis) and
$\psi_n(t) = \sin((n+1)\pi t/L)/\sqrt{L}$ (odd basis). Define the
\emph{difference matrix} $D_N(r)$ by
\begin{equation}
\label{eq:DN}
D_{nm}(r) = \langle\varphi_n, (S_{rL}+S_{-rL})\varphi_m\rangle - \langle\psi_n, (S_{rL}+S_{-rL})\psi_m\rangle.
\end{equation}

\begin{theorem}[Shift Parity Lemma, \cite{Geiger2026partII}]
\label{thm:SPL}
For every $N \ge 3$ and every $r \in (0, 2)$, $\lambda_{\min}(D_N(r)) < 0$.
\end{theorem}

This is a purely algebraic, closed-form statement about trigonometric
matrices; the entries of $D_N(r)$ are explicit sums of $\sin$ and $\cos$
with linear and constant coefficients in $r$.

\subsection{Frontier dominance}

The cumulative even--odd gap $\Delta(\lambda) = \lambda_1^+(\lambda) - \lambda_1^-(\lambda)$
is driven by primes in the \emph{frontier band} $[\lambda/e, \lambda]$: their
count is $\pi(\lambda) - \pi(\lambda/e) \sim \lambda/\log\lambda$ by the PNT,
and they are the ones for which $\log p / L \in [1-1/\log\lambda, 1]$. By
Theorem~\ref{thm:SPL}, each such prime contributes a negative summand of
order $-c\,\log p / \sqrt{p}$ to the difference, with the sum dominating any
archimedean or small-prime correction. Numerically, in \cite{Geiger2026partII}
frontier primes are shown to contribute $> 99\%$ of the gap for all tested
$\lambda \le 1.3\cdot 10^6$.

\subsection{NE-A and NE-B}

Define the \emph{prime shift multiplier}
\begin{equation}
\label{eq:Mxi}
M(\xi) = -2\,\mathrm{Re}\Big[\frac{\zeta'}{\zeta}\Big(\tfrac{1}{2}+i\xi\Big)\Big].
\end{equation}

\begin{theorem}[NE-A, \cite{Geiger2026partII}]
\label{thm:NEA-riemann}
$\{\xi \in \RR : M(\xi) < 0\}$ has positive Lebesgue measure. In particular,
the prime shift operator is not $\mathrm{PF}_\infty$-total positive.
\end{theorem}

\begin{theorem}[NE-B, \cite{Geiger2026partII}]
\label{thm:NEB-riemann}
For $N \in \{3,5,7,10,15\}$ and any sufficiently dense sample of $r \in (0,2)$,
the only symmetric matrix $T \in \mathrm{Sym}_N(\RR)$ satisfying
$[T, D_N(r_i)] = 0$ for all $i$ is $T = cI$. Consequently no nontrivial
universal commuting operator exists at these truncations; the full-space
conjecture \cite[Conjecture~5.3]{Geiger2026partII} asserts the same for
all $N$.
\end{theorem}

NE-B is the technically decisive structural result: it expresses that the
primes have no hidden common symmetry. The conceptual content is
multiplicative independence.


% =============================================================================
\section{The Selberg Setting}
\label{sec:selberg-setting}
% =============================================================================

We briefly record the standard setup.

\subsection{Compact hyperbolic surfaces}

Let $M = \Gamma\backslash\HH$, where $\HH = \{x+iy : y > 0\}$ is the upper
half-plane with metric $ds^2 = (dx^2+dy^2)/y^2$ and $\Gamma \subset \mathrm{PSL}(2,\RR)$
is a cocompact discrete torsion-free subgroup. Then $M$ is a compact
hyperbolic surface of genus $g \ge 2$. The Laplace--Beltrami operator
$\Delop$ acts on $L^2(M)$ as a non-negative self-adjoint operator with
purely discrete spectrum
\begin{equation}
0 = \lambda_0 < \lambda_1 \le \lambda_2 \le \dots,\qquad \lambda_j \to \infty.
\end{equation}
Writing $\lambda_j = 1/4 + r_j^2$ with $r_j \in \RR \cup i[-1/2, 1/2]$, the
spectral parameters $r_j$ encode the quantum spectrum.

\subsection{Primitive closed geodesics}

The conjugacy classes of hyperbolic elements in $\Gamma$ correspond bijectively
to (unoriented) closed geodesics on $M$. A geodesic is \emph{primitive} if it
is not a positive iterate of another. Denote by $\mathcal{P}_\Gamma$ the set
of primitive closed geodesics, and for $\gamma \in \mathcal{P}_\Gamma$ by
$\ell_\gamma$ its length.

The Prime Geodesic Theorem \cite{Huber1959, Sarnak1981} states
\begin{equation}
\label{eq:PGT}
\pi_\Gamma(T) := \#\{\gamma \in \mathcal{P}_\Gamma : \ell_\gamma \le T\} \sim \frac{e^T}{T},\qquad T \to \infty.
\end{equation}
This is the geodesic analogue of the PNT. Crucially, \emph{the density is
exponential, not polynomial}: $\pi_\Gamma(T)$ grows like $e^T/T$, whereas
$\pi(x)$ grows like $x/\log x = (e^{\log x})/\log x$. In other words, the
``prime-length-spectrum'' of geodesics, seen on the $\log$-axis, has density
$e^T/T$ (in the variable $T = \ell_\gamma$), comparable to the density of
primes on the $\log$-axis (in the variable $u = \log p$) only after the
change of variable.

\subsection{Selberg zeta function}

\begin{definition}
\label{def:Selberg-zeta}
The Selberg zeta function is
\begin{equation}
\label{eq:Sel-def}
\Sel(s) = \prod_{\gamma \in \mathcal{P}_\Gamma}\prod_{k=0}^\infty \bigl(1 - e^{-(s+k)\ell_\gamma}\bigr),\qquad \mathrm{Re}(s) > 1.
\end{equation}
\end{definition}

$\Sel(s)$ extends meromorphically to $\CC$, with:

\begin{itemize}
\item \emph{Trivial zeros} at $s = -k$, $k \in \NN_0$ (and further trivial zeros
associated with topology; see \cite{Hejhal1976}).
\item \emph{Non-trivial zeros} at $s = 1/2 \pm ir_j$ for each
Laplace eigenvalue $\lambda_j = 1/4 + r_j^2 > 0$.
\item A zero at $s = 1$ (from $\lambda_0 = 0$).
\end{itemize}

\begin{theorem}[Selberg spectral zero description and strict-line criterion, \cite{Selberg1956}]
\label{thm:Sel-RH}
The non-trivial spectral zeros of $\Sel(s)$ attached to Laplace eigenvalues
$\lambda_j=1/4+r_j^2$ occur at $s=1/2\pm ir_j$. Hence eigenvalues
$\lambda_j\ge 1/4$ give zeros on the line $\mathrm{Re}(s)=1/2$, whereas
exceptional eigenvalues $0<\lambda_j<1/4$ give real zeros
$s=1/2\pm\delta_j$ in the critical strip. In particular, the strict
critical-line statement holds under the spectral gap condition
$\lambda_1(\Delop)\ge 1/4$.
\end{theorem}

This follows from the non-negativity of $\Delop$: $\lambda_j \ge 0$ implies
$r_j \in \RR \cup i[-1/2, 1/2]$.  Non-negative real $r_j$ (i.e.\ $\lambda_j \ge
1/4$) give zeros on the critical line $\mathrm{Re}(s)=1/2$.  Exceptional
eigenvalues $\lambda_j \in (0,1/4)$, however, produce $r_j = i\delta_j$ with
$\delta_j \in (0,1/2)$, and thereby yield zeros at $s = 1/2 \pm \delta_j$
that are \emph{not} on the critical line.  Selberg-RH in the strict spectral
sense (all non-trivial spectral zeros in the critical strip on
$\mathrm{Re}(s)=1/2$) therefore requires the
spectral gap condition $\lambda_1(\Delop) \ge 1/4$.  This is Selberg's
eigenvalue conjecture \cite{Selberg1965}: it is established in the weaker form
$\lambda_1 \ge 3/16$ for congruence subgroups but remains open as
$\lambda_1 \ge 1/4$ in full generality.  The v2.0 analysis in
Section~\ref{sec:v2-to-selberg} is carried out under this hypothesis.

\begin{remark}[Functional equation]
\label{rem:func-eq}
The completed Selberg zeta function satisfies the functional equation
$\Sel(s) = \Sel(1-s) \cdot \Xi_\Gamma(s)$, where $\Xi_\Gamma(s)$ is an
explicit meromorphic factor depending on $\mathrm{Area}(M)$ and the genus of
$M$ \cite{Hejhal1976}.  In particular, non-trivial zeros occur in pairs
$\rho$, $1-\rho$, confirming their symmetry about $\mathrm{Re}(s)=\tfrac{1}{2}$.
\end{remark}

\subsection{Selberg explicit formula}

The Selberg trace formula \cite{Selberg1956, Hejhal1976} provides the
geodesic analogue of the Weil explicit formula. For a pair
$(h, \hat h)$ of sufficiently rapidly decaying functions related by Fourier
transform, the trace formula reads
\begin{equation}
\label{eq:Selberg-trace}
\sum_{j=0}^\infty h(r_j) = \frac{\mathrm{Area}(M)}{4\pi}\int_\RR h(r)\,r\tanh(\pi r)\,dr
+ \sum_{\gamma \in \mathcal{P}_\Gamma}\sum_{m=1}^\infty \frac{\ell_\gamma\,\hat h(m\ell_\gamma)}{2\sinh(m\ell_\gamma/2)}.
\end{equation}
The left-hand side is the \emph{spectral side} (a sum over eigenvalues of
$\Delop$); the right-hand side is the \emph{geometric side} (an identity term
plus a sum over geodesics). This plays exactly the role of
\begin{equation*}
\sum_\rho \Phi(\rho) = \text{(archimedean term)} - \sum_p\sum_{m\ge 1}\frac{\log p}{p^{m/2}}(\hat\Phi(m\log p) + \hat\Phi(-m\log p))
\end{equation*}
in the Weil explicit formula: primes $\leftrightarrow$ primitive geodesics,
$\log p \leftrightarrow \ell_\gamma$, non-trivial zeros $\rho \leftrightarrow$
spectral parameters $r_j$.


% =============================================================================
\section{Applying v2.0 to Selberg}
\label{sec:v2-to-selberg}
% =============================================================================

\subsection{Selberg geodesic shifts and Shift Parity Lemma}

Set $L = \log\lambda$. On $L^2([-L, L])$, define the Selberg prime shift operator
\begin{equation}
\label{eq:PSel}
(P^\Gamma_\lambda f)(t) = \sum_{\gamma \in \mathcal{P}_\Gamma, \ell_\gamma \le 2L}\sum_{m=1}^{\lfloor 2L/\ell_\gamma\rfloor}
\frac{\ell_\gamma}{2\sinh(m\ell_\gamma/2)}\bigl[f(t-m\ell_\gamma) + f(t+m\ell_\gamma)\bigr]\mathbf{1}_{[-L,L]}.
\end{equation}
In the cosine / sine basis of Section~\ref{sec:v2-recollect}, define the
\emph{Selberg difference matrix}
\begin{equation}
\label{eq:DNSel}
D^\Gamma_{nm}(r) = \langle\varphi_n, (S_{rL}+S_{-rL})\varphi_m\rangle - \langle\psi_n, (S_{rL}+S_{-rL})\psi_m\rangle.
\end{equation}
This is \emph{the same matrix} $D_N(r)$ as in the Riemann setting: it depends
only on the normalized shift $r$, not on the source of the shift (prime length
or geodesic length). Consequently:

\begin{proposition}[Transfer of Shift Parity]
\label{prop:SPL-transfer}
The Shift Parity Lemma (Theorem~\ref{thm:SPL}) holds verbatim in the Selberg
setting: for every $N \ge 3$, every $r \in (0,2)$,
$\lambda_{\min}(D^\Gamma_N(r)) < 0$.
\end{proposition}

\begin{proof}
$D^\Gamma_N(r) = D_N(r)$ as matrices. The Shift Parity Lemma depends only on
the harmonic structure of the cosine and sine bases on $[-L, L]$, and on the
fact that $S_{rL}+S_{-rL}$ is the sum of two shifts by the same amount in
opposite directions. The number-theoretic or geometric origin of $r$ is
immaterial.
\end{proof}

\begin{remark}
This is the algebraic universality. In the v2.0 framework for the Riemann
zeta, the Shift Parity Lemma is the bedrock on which everything rests. It
transfers to \emph{every} setting where zeros are parametrized by an explicit
formula with real, positive lengths playing the role of $\log p$. This
includes Dirichlet L-functions, Dedekind zeta functions, Selberg zeta,
Ruelle zeta, and in principle any Hecke or automorphic L-function.
\end{remark}

\subsection{Selberg Weil quadratic form}

On $L^2([-L, L])$, define the \emph{Selberg--Weil quadratic form}
\begin{equation}
\label{eq:QW-Sel}
\QW^\Gamma_\lambda(f) = \frac{\mathrm{Area}(M)}{4\pi}K_\mathrm{arch}^\Gamma(f,f)
+ \sum_{\gamma \in \mathcal{P}_\Gamma, \ell_\gamma\le 2L}\sum_{m=1}^\infty
\frac{\ell_\gamma}{2\sinh(m\ell_\gamma/2)}\langle T_{m\ell_\gamma} f + T_{-m\ell_\gamma}f, f\rangle,
\end{equation}
\begin{remark}[Truncation is exact]
\label{rem:truncation-exact}
For $f \in L^2([-L,L])$ and $|u| > 2L$, the shift inner product vanishes:
$\langle T_u f + T_{-u}f,\, f\rangle = 0$, because $f$ is supported on $[-L,L]$
while $T_{\pm u}f$ has support shifted outside $[-L,L]$.  Consequently
every geodesic iterate with $m\ell_\gamma > 2L$ contributes \emph{exactly} zero
to $\QW^\Gamma_\lambda$, with no approximation error.  The sum in
\eqref{eq:QW-Sel} is therefore not an approximation but an exact expression
for the quadratic form on $L^2([-L,L])$.
\end{remark}

where $K_\mathrm{arch}^\Gamma$ is the archimedean analogue
\begin{equation}
\label{eq:arch-Sel}
K_\mathrm{arch}^\Gamma(f,f) = \int_0^\infty[\langle T_u f + T_{-u}f, f\rangle - 2\,c_\Gamma^{\mathrm{id}}(u)\|f\|^2]\cdot \kappa^\Gamma(u)\,du,
\end{equation}
obtained by translating the identity term in \eqref{eq:Selberg-trace} into
shift operators.
\begin{remark}[Convergence and renormalisation at $u=0$]
\label{rem:arch-renorm}
Equation~\eqref{eq:arch-Sel} is the $u$-domain expression of the identity
term in the Selberg trace formula applied to $h(r) := |\hat f(r)|^2$; for
$f$ in a dense subspace (e.g.\ $C_c^\infty([-L,L])$), both sides converge
absolutely.  The pointwise singularity of $\kappa^\Gamma$ at $u=0$ is
matched by that of $c_\Gamma^{\mathrm{id}}$, and the bracket is the
regularised identity contribution; cf.\ \cite[Ch.~6]{Hejhal1976}.
\end{remark}
The coefficient
\begin{equation}
\label{eq:cGamma-def}
c_\Gamma^{\mathrm{id}}(u) \;:=\; \frac{\mathrm{Area}(M)}{4\pi}
\int_0^\infty r\tanh(\pi r)\cos(ru)\,dr
\end{equation}
is the Fourier--cosine transform of the Plancherel measure density
$r\tanh(\pi r)$, encoding the identity-coset contribution from
\eqref{eq:Selberg-trace}; the integral is interpreted in the
tempered-distribution sense since $r\tanh(\pi r)$ is polynomially
bounded (cf.\ \cite[Ch.~6]{Hejhal1976}).  Here
\begin{equation}
\label{eq:kappa-Sel}
\kappa^\Gamma(u) := \frac{1}{4\pi}\,\frac{d}{du}\!\Big[\frac{1}{\sinh(u/2)}\Big],\qquad u > 0,
\end{equation}
is the Plancherel kernel of the Selberg transform \cite[eq.~(1.40)]{Iwaniec2002}.
Note that $\kappa^\Gamma(u) < 0$ for all $u > 0$, since
$\tfrac{d}{du}[\sinh(u/2)^{-1}] = -\tfrac{\cosh(u/2)}{2\sinh^2(u/2)} < 0$.
This is consistent with the positivity structure of $K_\mathrm{arch}^\Gamma$:
the bracket in \eqref{eq:arch-Sel} is negative when $c_\Gamma^{\mathrm{id}}(u)$ normalises
the identity operator, and the product of two negative quantities is positive.

The quadratic form $\QW^\Gamma_\lambda$ admits a parity decomposition. Its
minimum eigenvalue is attained by either an even or an odd eigenfunction.

\begin{definition}[Selberg even dominance]
\label{def:Sel-ED}
We say that \emph{even dominance} holds for $\QW^\Gamma_\lambda$ at parameter
$\lambda$ if the minimum eigenvalue of $\QW^\Gamma_\lambda$ is simple and
attained by an even eigenfunction.
\end{definition}

\begin{remark}[Trace formula as bridge between spectral and geometric spaces]
\label{rem:two-spaces}
The form $\QW^\Gamma_\lambda$ is defined entirely in the spectral parameter
space $L^2([-L,L])$, where $T_u$ denotes the truncated shift $f(\cdot + u)$
restricted to $[-L,L]$.  A geometric realization via the geodesic flow on
$L^2(SM)$ also exists: both descriptions yield the same spectral sum
$\sum_j \hat{h}(r_j)$ when evaluated against test functions in the domain of
the Selberg transform.  The Selberg trace formula \eqref{eq:Selberg-trace}
serves as the implicit bridge between the two spaces.  An explicit intertwiner
via the Selberg--Helgason transform on $L^2(SM)$ is deferred to future work.
This is a limitation of the present paper: we use the trace formula as an
equality of spectral distributions, not as a constructed unitary equivalence
between the geodesic-flow representation and the truncated interval shift
model.  All NE-B language below is therefore restricted to the spherical
trace-formula channel unless such an explicit intertwiner is supplied.
\end{remark}

\subsection{Frontier dominance in the exponential density regime}

The frontier band for Selberg is the set of geodesics with
$\ell_\gamma / L$ close to $1$, i.e.\ $\ell_\gamma \in [L-1, L]$. By the Prime
Geodesic Theorem \eqref{eq:PGT},
\begin{equation}
\label{eq:frontier-count-Sel}
\#\{\gamma : \ell_\gamma \in [L-1, L]\} = \pi_\Gamma(L) - \pi_\Gamma(L-1)
\sim \frac{e^L}{L}(1 - e^{-1}) = \frac{\lambda}{\log\lambda}(1 - e^{-1}).
\end{equation}
This is \emph{the same asymptotic order} as the Riemann frontier count
$\pi(\lambda) - \pi(\lambda/e)$, up to the constant $(1-e^{-1}) \approx 0.632$.

The frontier geodesic weight is, for $\ell_\gamma \sim L$,
\begin{equation}
\label{eq:frontier-weight-Sel}
w^\Gamma_\gamma := \frac{\ell_\gamma}{2\sinh(\ell_\gamma/2)} \sim \frac{L}{e^{L/2}} = \log\lambda \cdot \lambda^{-1/2}.
\end{equation}
Compare this with the Riemann frontier prime weight
$w_p^{\mathrm{Riem}} = \log p / p^{1/2} \sim L \cdot \lambda^{-1/2}$ for
$p \sim \lambda$. The weights agree asymptotically in order of magnitude.

Combining \eqref{eq:frontier-count-Sel} and the weight estimate with the
quantitative Shift Parity magnitude bound and the Selberg off-frontier /
archimedean remainder estimates gives the following conditional frontier
estimate:
\begin{equation}
\label{eq:frontier-bound-Sel}
\Delta^\Gamma(\lambda) := \lambda_1^+(\QW^\Gamma_\lambda) - \lambda_1^-(\QW^\Gamma_\lambda)
\le -c_\Gamma\sqrt{\lambda} + o(\sqrt{\lambda}),
\end{equation}
for a positive constant $c_\Gamma > 0$ depending on the surface $M$ (via
$\mathrm{Area}(M)$ and the Selberg zeta function near $s=1$).

\begin{lemma}[Conditional frontier dominance constant]
\label{lem:frontier-constant}
Assume that the Part-II quantitative SPL magnitude bound and the corresponding
off-frontier/archimedean remainder estimate hold for the Selberg kernel. For
every compact hyperbolic surface $M$ there exists
$\lambda_0 = \lambda_0(M) > 0$ such that the frontier dominance constant in
\eqref{eq:frontier-bound-Sel} satisfies
\begin{equation}
\label{eq:cGamma-explicit}
c_\Gamma \;\ge\; \tfrac{1}{2}(1-e^{-1}) \;>\; 0
\qquad\text{for all }\lambda \ge \lambda_0.
\end{equation}
\end{lemma}

\begin{remark}
The normalised unsigned frontier sum $\lambda^{-1/2}\sum_{\ell_\gamma\in[L-1,L]}
w^\Gamma_\gamma \to (1-e^{-1}) \approx 0.632$ as $\lambda\to\infty$.
The factor $\tfrac{1}{2}$ in \eqref{eq:cGamma-explicit} is conservative: it uses
the quantitative lower bound of \cite[Lem.~A.3]{Geiger2026partII}, which gives
signed magnitude $\ge\tfrac{1}{2}w^\Gamma_\gamma$ per frontier geodesic.
Proposition~\ref{prop:SPL-transfer} provides the sign ($\lambda_{\min}<0$)
alone; the magnitude bound is the content of the Part-II lemma.
Under the sharp bound, $c_\Gamma \to (1-e^{-1})$ asymptotically.
\end{remark}

\begin{proof}[Proof sketch]
By \eqref{eq:frontier-count-Sel}, the frontier band $[L-1,L]$ contains
\[
(1-e^{-1})\,\frac{e^L}{L}\,(1+o(1))
=
(1-e^{-1})\,\frac{\lambda}{\log\lambda}\,(1+o(1))
\]
primitive geodesics. Each contributes weight
$w^\Gamma_\gamma = L/(2\sinh(L/2)) \sim \log\lambda\cdot\lambda^{-1/2}$ by
\eqref{eq:frontier-weight-Sel}. By \cite[Lem.~A.3]{Geiger2026partII},
each frontier summand is negative with magnitude at least
$\tfrac{1}{2}w^\Gamma_\gamma$ for all large $\lambda$.  Summing over the frontier
band and dividing by $\sqrt\lambda$ gives
\[
\sum_{\ell_\gamma\in[L-1,L]} w^\Gamma_\gamma \;\ge\;
(1-e^{-1})\,\frac{\lambda}{\log\lambda}\cdot
\frac{\log\lambda}{\sqrt\lambda}\cdot\tfrac{1}{2}\,(1+o(1))
= \tfrac{1}{2}(1-e^{-1})\sqrt\lambda\,(1+o(1)),
\]
yielding $c_\Gamma \ge (1-e^{-1})/2$.
\end{proof}

\begin{remark}
The factor $(1-e^{-1}) \approx 0.632$ is the geometric ratio of frontier
geodesic density (exponential regime) to frontier prime density
(polynomial regime): Selberg has $(1-e^{-1})\lambda/\log\lambda$ frontier
geodesics where Riemann has $\lambda/\log\lambda$ frontier primes.
This accounts for the only structural difference between the two frontier
estimates.
\end{remark}

\subsection{Conditional reduction via v2.0}

\begin{theorem}[v2.0 Selberg-RH, conditional]
\label{thm:v2-Sel-RH}
Let $M = \Gamma\backslash\HH$ be a compact hyperbolic surface satisfying
$\lambda_1(\Delop) \ge 1/4$.  Assume the following three C2 conditions for
$\QW^\Gamma_\lambda$:
\begin{enumerate}[label=(C2\alph*)]
\item $\QW^\Gamma_\lambda$ has a lower-bounded self-adjoint realization on the
stated interval model;
\item for all $\lambda \ge \lambda_0(\Gamma)$ its lowest eigenvalue is simple
and attained by an even minimizer $\eta_{\lambda,\Gamma}$;
\item the Fourier/Selberg transform of $\eta_{\lambda,\Gamma}$ is identified,
in the relevant exact model or limit, with the completed Selberg spectral
factor whose zeros encode $\Sel$.
\end{enumerate}
Then all non-trivial spectral zeros of $\Sel(s)$ in the critical strip lie on
$\mathrm{Re}(s) = 1/2$.  Without (C2c), Connes' mechanism only gives a
statement about the zeros of the minimizer transform, not about the spectral
zeros of $\Sel$.
\end{theorem}

For readability, the C2 gate is recorded separately from the theorem
statement.  Table~\ref{tab:c2-minimizer-ledger} distinguishes what is already
used inside this paper from the operational ledger that remains a companion-
level audit rather than a theorem-level closure.  The no-exceptional-eigenvalue
assumption is shown as a separate SG/NoExc gate; it is not part of C2b.

\begin{table}[ht]
\centering
\small
\setlength{\tabcolsep}{3pt}
{\renewcommand{\arraystretch}{1.16}
\begin{tabularx}{\textwidth}{@{}L{0.18\textwidth}YYY@{}}
\toprule
Gate & Object recorded & Current use in this paper & Failure or revision signal \\
\midrule
SPH: spherical channel
  & Selberg transform and radial \(K\)-bi-invariant convolution acting through
    \(\widehat{k}(r_j)\), with Casimir/Laplace commutation.
  & Positive-control context for NE-B failure; it is not itself a C2 gate and
    does not identify individual geodesic-flow pullbacks with interval shifts.
  & A non-radial or full \(SM\)-shift model is used as if it were the truncated
    \(L^2([-L,L])\) model. \\
C2a: interval realization
  & Lower-bounded self-adjoint interval model for \(\QW^\Gamma_\lambda\).
  & Explicit v2.0 hypothesis for applying the Connes mechanism.
  & The interval model is only asserted by analogy or by the spherical channel,
    without a stated self-adjoint realization. \\
C2b: parity/minimizer gate
  & Simple lowest eigenvalue and even minimizer for the interval realization.
  & Explicit v2.0 hypothesis; supported by the Selberg/Casimir channel but not
    derived here from the spectral gap.
  & The minimum is degenerate, odd, or selected only after observing the target
    Selberg factor. \\
SG/NoExc: spectral-line input
  & \(\lambda_1(\Delop)\ge 1/4\) and absence of exceptional eigenvalues.
  & Explicit hypothesis for the strict critical-line form; the spectral zero
    description itself remains classical.
  & An exceptional eigenvalue \(0<\lambda_j<1/4\) produces real zeros
    \(1/2\pm\delta_j\). \\
C2c: minimizer/window identification
  & Target window \(I\), target source, center \(\nu\), outer gap, residual,
    off-window mass, and negative control.
  & Open ledger; it must be filled before the v2.0 minimizer is identified with a
    Selberg spectral factor.
  & The candidate has the right Rayleigh center but wrong tail or off-window mass
    (\texttt{same\_rayleigh\_}\allowbreak\texttt{wrong\_tail}). \\
Decision layer
  & \texttt{pass}, \texttt{ambiguous}, or \texttt{fail} for each surface/window
    candidate.
  & Positive-control audit, not a proof upgrade; it localizes what remains to be
    checked for JST or companion work.
  & Missing table, same-run target choice, or absent negative controls. \\
\bottomrule
\end{tabularx}}
\caption{C2/minimizer-identification ledger for the Selberg positive-control reading.}
\label{tab:c2-minimizer-ledger}
\end{table}

\begin{proof}[Proof sketch]
The structure of the argument is identical to the v2.0 proof for Riemann.

\smallskip\noindent
\textbf{Step 1 (pointwise even preference, including iterates).} By
Proposition~\ref{prop:SPL-transfer}, for every $\gamma \in \mathcal{P}_\Gamma$
and every $m \ge 1$ with $m\ell_\gamma \le 2L$, the difference matrix
$D^\Gamma_N(m\ell_\gamma / L)$ has $\lambda_{\min} < 0$.  Here $m > 1$
corresponds to the $m$-fold iterate of $\gamma$; the SPL applies to iterates
on exactly the same footing as to primitive geodesics, since the shift length
$m\ell_\gamma$ enters the SPL as a scalar parameter $r = m\ell_\gamma/L$.
By Remark~\ref{rem:truncation-exact}, iterates with $m\ell_\gamma > 2L$
contribute exactly zero and need not be handled separately.  Hence each individual
geodesic summand (primitive or iterated), weighted by the positive coefficient
$\ell_\gamma/(2\sinh(m\ell_\gamma/2))$, contributes negatively to the
even-minus-odd difference.

\smallskip\noindent
\textbf{Step 2 (frontier dominance).} By \eqref{eq:frontier-count-Sel} and
\eqref{eq:frontier-bound-Sel}, the aggregate contribution of frontier geodesics
to $\Delta^\Gamma(\lambda)$ is negative of order $-c_\Gamma \sqrt{\lambda}$,
while the combined contribution of non-frontier geodesics plus the archimedean
term is controlled by the Selberg analogue of the Part-II remainder estimates.
This is exactly the conditional estimate recorded in
Lemma~\ref{lem:frontier-constant}.

\smallskip\noindent
\textbf{Step 3 (C2b parity/minimizer gate).} Simplicity and evenness of the minimum are not
proved here from generic simplicity or from a direct identification with
$\lambda_1(\Delop)$.  They are precisely condition (C2b).  The Selberg/Casimir
channel is strong evidence for the right parity structure, but the present
paper treats this as a stated gate.  C2b is not the separate SG/NoExc
spectral-gap gate.

\smallskip\noindent
\textbf{Step 4 (Connes' reduction).} Theorem~\ref{thm:connes} is the
\emph{algebraic mechanism}: given ingredients~(i) a lower-bounded self-adjoint
operator on a symmetric interval and (ii) a simple minimum with even eigenfunction,
it forces all zeros of the minimum eigenvector's Fourier transform onto the real
line.  Verifying ingredients~(i) and~(ii) is the $C2^\Gamma$ substep, while
identifying the resulting minimizer transform with the completed Selberg
factor is condition (C2c).  (In the three-component
decomposition $\mathrm{RH}_X = \mathrm{GBC}_X + C2_X + \mathrm{Hurwitz}_X$
developed in the companion work \cite{GeigerMeta2026}, in preparation, these
correspond to the Lie-origin branch plus the explicit minimizer-identification
gate.) For Riemann,
ingredient~(ii) is non-trivial; it is the content of v2.1
Proposition~A6 \cite{Geiger2026partII}. Applying Connes reduction with
$\QW^\Gamma_\lambda$ in place of $\QW_\lambda$, the zeros of
$\widehat{\eta}_{\lambda,\Gamma}$ are real.  Condition (C2c) is the additional
gate that transfers this real-zero statement to the completed Selberg factor.
Together with $\lambda_1(\Delop)\ge 1/4$, this gives the strict statement for
the non-trivial Selberg spectral zeros in the critical strip under the
theorem's explicit hypotheses.
\end{proof}

\begin{remark}[Conditionality and consistency]
Theorem~\ref{thm:v2-Sel-RH} is a conditional reduction: \emph{given} the C2
interval/minimizer gates and the separate no-exceptional-eigenvalue condition, the
strict spectral-zero statement follows by the v2.0 argument.  Since the classical
Selberg theory already gives the spectral zero description and ties the
strict line statement to the gap $\lambda_1(\Delop)\ge 1/4$, it verifies the
framework's conclusion independently at the correct hypothesis level.  The
value of the conditional reduction lies in \emph{stress-testing the reach of
the v2.0 method}: it applies in a case where a Hilbert--P\'olya operator
already exists, and the two routes reach the same conclusion by independent
means.  This agreement is evidence of consistency, not a proof upgrade.
\end{remark}

\begin{remark}[Even dominance as structural expectation]
\label{rem:even-dominance-structural}
Condition (C2b) in Theorem~\ref{thm:v2-Sel-RH} is a
\emph{structural expectation}: it is suggested by the positivity of the
hyperbolic geodesic weights and by the Casimir/spherical trace-formula channel,
but it is not proved here.  A fully rigorous proof would require the same
frontier-dominance estimates as in the Riemann Part~II
\cite{Geiger2026partII}, adapted to the exponential geodesic density, and the
minimizer-identification gate (C2c). Lemma~\ref{lem:frontier-constant}
provides only the conditional frontier input. For Selberg, the structural
expectation is additionally supported by the
independent operator-theoretic proof (Theorem~\ref{thm:Sel-RH}), which
confirms the conclusion without using even dominance at all.
\end{remark}

\subsection{NE-A for Selberg}

Define the Selberg prime shift multiplier
\begin{equation}
\label{eq:Mxi-Sel}
M^\Gamma(\xi) = -2\,\mathrm{Re}\Big[\frac{\Sel'}{\Sel}\Big(\tfrac{1}{2}+i\xi\Big)\Big].
\end{equation}

\begin{proposition}[NE-A for Selberg]
\label{prop:NEA-Sel}
The set $\{\xi \in \RR : M^\Gamma(\xi) < 0\}$ has positive Lebesgue measure.
\end{proposition}

\begin{proof}[Proof sketch]
By the Hadamard product representation of the Selberg zeta function on a
compact hyperbolic surface (\cite[Ch.~10]{Iwaniec2002}), one has
\begin{equation}
\label{eq:hadamard-Sel}
\begin{aligned}
\frac{\Sel'}{\Sel}(s)
&= b_\Gamma + \sum_{\rho}\Big[\frac{1}{s-\rho} + \frac{1}{\rho}\Big] \\
&\quad + (\text{contributions from trivial zeros and the topological factor}),
\end{aligned}
\end{equation}
where $\rho$ runs over the non-trivial zeros (on $\mathrm{Re}(s)=1/2$ under
Theorem~\ref{thm:Sel-RH}) and $b_\Gamma \in \RR$ is a constant depending on~$\Gamma$.

Setting $s = 1/2 + i\xi$ and taking real parts, each non-trivial zero
$\rho_j = 1/2 + ir_j$ contributes
\begin{equation}
\label{eq:NEA-re}
\mathrm{Re}\Big[\frac{1}{i(\xi-r_j)}+\frac{1}{1/2+ir_j}\Big]
  = \frac{1/2}{1/4+r_j^2} > 0.
\end{equation}
Thus the non-trivial zeros produce a \emph{positive constant} contribution
to $\mathrm{Re}[\Sel'/\Sel(1/2+i\xi)]$; they do not create sign oscillations
in~$M^\Gamma(\xi)$.

The sign oscillations arise from the trivial zeros (at $s = -k$, $k \ge 0$)
and the topological factor, whose contributions to $\mathrm{Re}[\Sel'/\Sel]$
are of the form $\sum_{k\ge 0} c_k/[(\xi^2+(k+1/2)^2)]$ with $c_k \in \RR$
of alternating sign. Combining with the explicit-formula identity term of the
Selberg trace formula \eqref{eq:Selberg-trace}, one finds by a direct
application with a suitable test function (localising in $\xi$) that
$M^\Gamma(\xi)$ takes both strictly positive and strictly negative values on
sets of positive Lebesgue measure. The argument is the direct analogue of
\cite[Lem.~A.3]{Geiger2026partII}; the only substitution is $\tanh(\pi r)$
in place of $\Gamma'/\Gamma(1/4+ir/2)$ in the archimedean kernel.
\end{proof}

NE-A therefore transfers. The prime shift multiplier is oscillatory in both
settings. In the Selberg case the sign changes come from the
trivial/topological/identity terms in the explicit formula, not from the
critical-line zeros themselves.


% =============================================================================
\section[NE-B Fails for Selberg: The Spherical Casimir Channel]{NE-B Fails for Selberg: The Spherical Casimir Channel}
\label{sec:ne-b-fails}
% =============================================================================

Here we come to the substantive difference between Selberg and Riemann.

\subsection{Statement}

\begin{proposition}[Casimir commutation in the spherical Selberg channel]
\label{thm:NEB-fail}
Let $G=\mathrm{PSL}(2,\RR)$ and $K=\mathrm{SO}(2)$, so
$M=\Gamma\backslash G/K$.  Let $R_k$ be the radial convolution operator on
$L^2(\Gamma\backslash G/K)$ associated with a $K$-bi-invariant test kernel
$k$.  Then, on the common invariant domain,
\begin{equation}
\label{eq:NEB-commut}
[\Delop, R_k]=0.
\end{equation}
Consequently the Selberg trace-formula channel has a non-trivial Casimir
commutant.  In this precise spherical sense, the Selberg analogue of the
NE-B obstruction is absent.
\end{proposition}

\subsection{Spherical radial convolution operators}

The objects that v2.0 manipulates directly are shifts on the
\emph{logarithmic interval} $[-L,L]$.  In the Selberg trace formula these
interval shifts are represented through test functions and their spherical
transforms.  The operatorial statement made here is therefore not a statement
about individual geodesic-flow pullbacks on $L^2(SM)$; it is a statement about
the spherical convolution algebra on $L^2(\Gamma\backslash G/K)$.

For a compactly supported $K$-bi-invariant kernel $k$ on $G$, set
\[
(R_k f)(x) := \int_G k(g) f(xg)\,dg
\]
on the automorphic quotient, whenever the integral is defined.  These radial
operators are the spectral side of the Selberg trace formula: under the
Selberg/Harish--Chandra transform they act on each Laplace eigenspace with
scalar multiplier $\hat{k}(r_j)$.

\subsection{Proof of Proposition~\ref{thm:NEB-fail}}

\begin{proof}
The Laplace--Beltrami operator $\Delop$ is the image of the Casimir element in
the center of $\mathcal{U}(\mathfrak{sl}_2)$.  Therefore it commutes with the
regular representation and, in particular, with every $K$-bi-invariant
convolution operator $R_k$ on the spherical quotient.  Equivalently, the
Selberg/Harish--Chandra transform diagonalizes both $\Delop$ and $R_k$: on a
Laplace eigenspace with spectral parameter $r_j$, $\Delop$ acts by
$1/4+r_j^2$ and $R_k$ acts by the scalar $\hat{k}(r_j)$.  Hence their
commutator vanishes on the common invariant domain.
\qedhere
\end{proof}

\begin{remark}[No individual phase-eigenvalue claim]
The paper does not use the stronger and generally unsafe assertion that an
individual geodesic-flow element satisfies
$T^\Gamma_t\psi_j=e^{itr_j}\psi_j$ for arbitrary Laplace eigenfunctions.
The safe statement is the spherical one: radial convolution operators are
diagonalized by the Selberg/Harish--Chandra transform with scalar multiplier
$\hat{k}(r_j)$.
\end{remark}

\begin{remark}[Why this is possible for Selberg but not for Riemann]
The argument depends on the existence of a spherical homogeneous channel
$\Gamma\backslash G/K$ and its central Casimir.  Primitive geodesics enter the
trace formula through conjugacy classes and radial kernels, while the Casimir
diagonalizes the same spectral side.  Primes, by contrast, do not come with a
known comparable spherical convolution algebra and central differential
operator.  This is what NE-B captures at the level used here: the tested
Riemann shift model has no corresponding universal commutant.
\end{remark}

\subsection{Consequence for the numerical commutator test}

In the spherical Selberg model, repeating the numerical commutator test of
\cite[Section~5.3]{Geiger2026partII} should yield a \emph{large} kernel for
the commutator constraint matrix $C_N$ -- in fact a kernel whose dimension
grows with the number of distinct Laplace eigenvalues sampled, because every
polynomial in $\Delop$ is a commutant. The singular-value spectrum would show
no ``isolated-$1$-dim-null''-structure; instead many small singular values
would appear, reflecting the infinite-dimensional commutant.

This is the diagnostic: \emph{if you run the v2.0 commutator test on a
genuine Hilbert--P\'olya system, it detects the operator}. In the Riemann
case, the tested truncation detects nothing nontrivial -- and this is, by the
logic of v2.0, evidence that this specific commutant route is absent.


% =============================================================================
\section{Hilbert--P\'olya as a Special Case of v2.0}
\label{sec:meta}
% =============================================================================

We now articulate the meta-principle that emerges from combining the
Riemann and Selberg analyses.

\subsection{The Hilbert--P\'olya programme revisited}

The classical Hilbert--P\'olya conjecture proposes to prove RH by constructing
a self-adjoint operator $H$ on a Hilbert space whose spectrum is the set of
imaginary parts $\{\gamma_k\}$ of the non-trivial Riemann zeros. Such an
operator, if it exists, would give RH as an immediate corollary
(self-adjointness $\Rightarrow$ real eigenvalues $\Rightarrow$ zeros on
the critical line).

One hundred years of effort have produced candidate operators (Berry--Keating
$H = xp$ \cite{BerryKeating1999}, Connes' adelic Dirac operator
\cite{Connes1999}, and others) but none that provably realize the Riemann
spectrum.

\subsection{The v2.0 alternative}

The v2.0 programme circumvents the Hilbert--P\'olya requirement by
establishing RH through a \emph{quadratic-form inequality} (even dominance
of $\QW_\lambda$), rather than through an operator-spectral identification.
The ingredients are:
\begin{itemize}
\item A parity decomposition of the ambient Hilbert space.
\item An algebraic \emph{pointwise} parity advantage (Shift Parity Lemma).
\item An analytic \emph{aggregation} mechanism (frontier dominance).
\item An argument that no simpler route applies (NE-A, NE-B).
\end{itemize}
Crucially, \emph{no operator is constructed}. The route is via the
quadratic form directly.

\subsection{Selberg as calibration: the spherical operatorial channel exists}

Theorem~\ref{thm:v2-Sel-RH} shows that v2.0 conditionally reduces the strict
statement for the non-trivial Selberg spectral zeros in the critical strip to the C2/minimizer gates, even dominance of
$\QW^\Gamma_\lambda$, and $\lambda_1(\Delop)\ge 1/4$.  In this case, the
spherical Casimir channel also exists (Proposition~\ref{thm:NEB-fail}), and
classical spectral theory independently confirms the conclusion.  The two
routes coexist:
\begin{itemize}
\item Selberg's original spectral theory (unconditional spectral
  parametrisation; strict critical line under $\lambda_1(\Delop)\ge 1/4$).
\item The v2.0 reduction (quadratic: the C2/minimizer gates and even dominance
  imply the strict spectral-zero statement under the stated gap hypothesis).
\end{itemize}
Both arrive at the same conclusion by different routes once their respective
hypotheses are supplied. The v2.0 method is therefore compatible with the
operator-based route; it does not by itself replace the C2/minimizer
identification.

\subsection{Riemann as the complementary case}

In the Riemann setting, NE-B (in the full-space conjectural form) asserts that
the tested universal-commutant route is absent. If this conjecture holds, then
the particular Hilbert--P\'olya strategy modeled by these shifts is blocked.
This is the \emph{scientifically interesting} case: v2.0 seeks a proof route
that does not start by postulating such a commutant.

\subsection{Meta-theorem}

\begin{theorem}[v2.0 meta-principle, informal]
\label{thm:meta}
Let $Z$ be a zeta-like function admitting an explicit formula of the
Riemann--Weil type (with real-valued logarithmic shifts and positive weights).
Then:
\begin{enumerate}[label=(\roman*)]
\item The v2.0 ingredients (Shift Parity Lemma, frontier dominance,
parity decomposition of the associated Weil quadratic form) transfer to $Z$
with modifications only in the density and weight constants.
\item The v2.0 route to the corresponding Riemann Hypothesis is a candidate
route whenever an appropriate frontier-density bound and C2/minimizer gate hold.
\item Failure of the relevant NE-B obstruction is evidence for an
operatorial/spherical Hilbert--P\'olya channel for $Z$.
\item When such a channel exists, v2.0 is a consistency check against a
classical operatorial route. When it does not, NE-B blocks this specific
commutant strategy, not every imaginable Hilbert--P\'olya realization.
\end{enumerate}
\end{theorem}

We formulate this as a \emph{structural principle} rather than a formal
theorem, because it depends on what one takes as the admissible class of
zeta-like functions. For classical L-functions satisfying a standard
functional equation and explicit formula (Selberg class, Connes' family),
these items are expected to hold only after the relevant operator class and
C2/minimizer-identification gates are specified.

\begin{remark}[Refinement in the companion meta-paper]
\label{rem:meta-paper-refinement}
The informal meta-principle of Theorem~\ref{thm:meta} has been made
technically precise in the companion paper \cite{GeigerMeta2026} as the
\emph{Meta-Theorem} (Thm.~thm:meta there), via the three-component
decomposition
\[
\mathrm{RH}_X \;=\; \mathrm{GBC}_X \;+\; C2_X \;+\; \mathrm{Hurwitz}_X.
\]
In that language, items~(iii)--(iv) of our informal Theorem~\ref{thm:meta}
are the statement that $C2_X$ is trivialised in the Lie-origin branch and
open (equivalent to classical Hilbert--P\'olya) in the arithmetic branch.
The finer point noted in \cite{GeigerMeta2026}
is that $C2_X$ itself decomposes into a \emph{parity substep}
$C2_X^{\mathrm{parity}}$ (simple ground state with even eigenvector) and an
\emph{eigenvector-realisation substep} $C2_X^{\mathrm{eigenvec}}$
(Hermite-prolate approximation convergence). In the current CORE status,
the Riemann $C2/MS2$ block is closed in the CCM/Fourier branch by the
Zookeeper three-lemma microcluster closure.  This does not change the
Selberg proof obligation: for Selberg, both C2 substeps are trivialised
simultaneously by the Laplace--Casimir.  The Selberg case is therefore a
positive control for the normal form, not a dependent corollary of the
Riemann closure.
\end{remark}

\subsection{Restatement: the Weil quadratic form is fundamental}

The structural insight is:

\medskip
\noindent\emph{The Weil quadratic form $\QW$ is the fundamental object; the
Hilbert--P\'olya operator is a realization that may or may not exist.}

\medskip
For one hundred years, the Riemann Hypothesis has been viewed through the
Hilbert--P\'olya lens: find the operator, prove RH. The v2.0 framework
inverts this view: prove even dominance, prove RH, and investigate (as a
separate question) whether a Hilbert--P\'olya operator happens to exist.
In the Selberg spherical channel it does; in the tested Riemann shift model no
comparable commutant is available. Both cases are captured by the same
underlying framework.


% =============================================================================
\section{Applications and Outlook}
\label{sec:applications}
% =============================================================================

\subsection{Generalized L-functions}

The v2.0 method applies in principle to any zeta-like function with a
Riemann--Weil-type explicit formula. We expect the following:

\begin{itemize}
\item \emph{Dirichlet L-functions} $L(s, \chi)$ for a Dirichlet character
$\chi$: the explicit formula is the prime summation weighted by $\chi(p)$.
The Shift Parity Lemma transfers (the character introduces a phase, not a
sign-change of the geometric parity structure). Frontier dominance transfers
(with PNT replaced by the PNT for $L$-functions, Landau--Siegel bounds).
NE-B: the tested commutant obstruction is expected to remain live unless the
character introduces a new operatorial channel. In the current Zeta-Zoo
cartography, Dirichlet is not closed merely by importing the untwisted Riemann
package: it remains the Weg-D / $\chi$-twisted CCM follow-up branch. A concrete
Dirichlet write-up would make the meta-framework's predictions operationally
testable --- it is the natural next step after the present Selberg consistency
check.

\item \emph{Dedekind zeta functions} $\zeta_K(s)$ for a number field $K$: the
explicit formula sums over prime ideals $\mathfrak{p}$ with weight
$(\log N\mathfrak{p}) \cdot N\mathfrak{p}^{-m/2}$. The Shift Parity and
frontier mechanisms transfer. NE-B: conjecturally no comparable spherical
channel exists for generic fields, while special extra-symmetry cases need
separate treatment.

\item \emph{Automorphic L-functions} (Hecke, $\mathrm{GL}(n)$-cuspidal):
the explicit formula transfers. A Hilbert--P\'olya channel is available when
the L-function has a spectral interpretation via a Laplacian or Casimir
channel (which happens when the automorphic form is itself of geometric origin
--- e.g.\ Maass forms on a hyperbolic surface). Otherwise the NE-B analogue
remains a separate obstruction to test.
\end{itemize}

\subsection{Ruelle zeta for Axiom-A flows}

For an Axiom-A flow on a compact Riemannian manifold (a hyperbolic dynamical
system), the Ruelle zeta function is defined by
\begin{equation}
\zeta_R(s) = \prod_\tau \frac{1}{1 - e^{-s\ell_\tau}},
\end{equation}
where the product is over primitive periodic orbits $\tau$ with minimal
period $\ell_\tau$.  Results such as \cite{Dolgopyat1998} and subsequent
Anosov-zeta work concern analytic continuation, resonances, and decay
mechanisms; they should not be read here as a general critical-line theorem.
The following item is therefore only an outlook problem.

\begin{conjecture}[v2.0 Ruelle test problem]
\label{conj:Ruelle}
For a hyperbolic flow with topological entropy $h$, it is worth testing
whether a v2.0-style Weil quadratic form $\QW^\mathrm{R}_\lambda$ admits
frontier dominance and a useful C2/minimizer gate. Any critical-line-type
statement for $\zeta_R$ would require additional resonance-specific input and
is not established by the present Selberg argument.
\end{conjecture}

The point of the conjecture is diagnostic: for generic Axiom-A flows, no
Selberg-style symmetric-space Casimir is available, so the NE-B question and
the C2 gate would have to be rebuilt from resonance theory rather than
imported from the Selberg case.

\subsection{Ihara zeta and the discrete--continuum axis}
\label{sec:ihara}

The Ihara zeta function $\zeta_\Gamma(s)$ of a finite connected $(q+1)$-regular
graph $\Gamma$ is defined by
\begin{equation}
\zeta_\Gamma(s)^{-1}
  = (1-s^2)^{r-1}\det\!\bigl(I - sA + qs^2 I\bigr),
\end{equation}
where $A$ is the adjacency matrix and $r$ is the cycle rank of $\Gamma$
\cite{Ihara1966,Hashimoto1989,Bass1992}.  The determinant formula
is sometimes called the Ihara--Bass formula; the proofs by Hashimoto (1989)
and Bass (1992) established it in full generality.  The Ihara~RH is the statement that all
non-trivial poles lie on the circle $|s| = q^{-1/2}$, equivalent to
$\Gamma$ being a Ramanujan graph.  It has been proved for explicit
Ramanujan families constructed from $\mathrm{PGL}(2)$-arithmetic
\cite{LPS1988}.

\medskip
In the v2.0 / Zookeeper framework, the Ihara case is \emph{discrete--operatorial}
(SGE-YES), in parallel with the \emph{continuous--operatorial} Selberg case.
The graph Laplacian $L = (q+1)I - A$ commutes with all automorphisms of $\Gamma$
and hence with the edge-adjacency transfer matrix $T$.  This makes $L$ the
discrete analogue of the hyperbolic Laplacian $\Delta_\Gamma$: a universal
commutant that renders NE-B vacuous for Ihara, for exactly the same structural
reason as in the Selberg case.

The three-way comparison axis in the Zookeeper framework is:

\begin{center}
\footnotesize
\setlength{\tabcolsep}{4pt}
\begin{tabularx}{\textwidth}{@{}L{2.2cm}L{2.1cm}Y L{1.8cm}@{}}
\toprule
\textbf{Zeta function} & \textbf{Type} & \textbf{Commutant} & \textbf{v2.0 class} \\
\midrule
Selberg $Z_\Gamma(s)$ & continuous, geometric & $\Delta_\Gamma$ (Laplace--Beltrami) & SGE-YES \\
Ihara $\zeta_\Gamma(s)$ & discrete, combinatorial & $L = (q{+}1)I - A$ (graph Laplacian) & SGE-YES \\
Riemann $\zeta(s)$ & arithmetic & none (NE-B holds) & SGE-NO \\
\bottomrule
\end{tabularx}
\end{center}

\medskip\noindent
In both SGE-YES cases the strict RH-type statement is governed by a classical
spectral bound (Selberg's no-exceptional-eigenvalue condition
$\lambda_1(\Delop)\ge 1/4$; Ramanujan's $|$eigenvalue$| \le 2\sqrt{q}$),
and the v2.0 framework provides a \emph{consistent conditional re-derivation} via even dominance
of the corresponding Weil quadratic form.  For the Riemann zeta, no such
commutant exists; the v2.0 method is the only available structural route to RH.

\subsection{Physical interpretation}

The Selberg setting corresponds physically to a quantum particle on a
hyperbolic surface (``quantum chaos''); the Laplace eigenvalues are the
energy levels. The v2.0 Shift Parity Lemma states that the parity of the
eigenfunctions is biased toward even functions at the threshold. This is a
statement about the parity-class distribution of quantum eigenstates.

For Axiom-A flows (billiards, chaotic cavities), the v2.0 framework predicts
GUE-type statistics for the Ruelle zeros, as in the BGS (Bohigas--Giannoni--Schmit)
conjecture \cite{BGS1984}; see also \cite{RudnickSarnak1996} for related GUE
statistics of zeros of principal $L$-functions. The v2.0 methodology may offer a
structural explanation of this conjecture.


% =============================================================================
\section{Zookeeper Consistency: Selberg as an SGE-YES Test}
\label{sec:zookeeper-selberg}

The Selberg case is also a useful calibration point for the post-Zookeeper
version of the Zeta Zoo.  In this language, Selberg is an SGE-YES system:
the relevant operator exists classically, and the v2.0 transfer must
therefore reproduce a known positive case rather than create a new
assertion from scratch.

Table~\ref{tab:transfer-slots} records the five transfer slots and
their occupants in the Selberg and Riemann settings.

\begin{table}[ht]
\centering
\footnotesize
\setlength{\tabcolsep}{5pt}
\caption{The five SGE-YES transfer slots for Selberg vs.\ Riemann.
  A tick (\checkmark) indicates the slot is classically occupied;
  ``v2.0'' indicates it is established by the v2.0 framework; ``hyp.''
  marks the extra spectral-gap hypothesis needed for the strict-line statement.}
\label{tab:transfer-slots}
\begin{tabularx}{\textwidth}{@{}p{2.4cm}YY@{}}
\toprule
\textbf{Slot} & \textbf{Selberg (SGE-YES)} & \textbf{Riemann} \\
\midrule
Operator / form
  & $\Delop$ (classical Casimir; commutes with spherical radial convolutions; interval-shift intertwiner not constructed) \checkmark
  & $\QW_\lambda$ (no classical HP operator; NE-B holds) \\[4pt]
Defect
  & Casimir residual $S_\Gamma(L,I)=r_\Gamma/g_I$ (companion proposal)
  & Gap to even dominance (conjectural $\Rightarrow$ established by v2.0) \\[4pt]
Approximation
  & Length cutoffs, truncated trace formula \checkmark
  & Galerkin truncation $D_N(r)$ (v2.0) \\[4pt]
Gap
  & No-exceptional gap $\lambda_1(\Delop)\ge \tfrac{1}{4}$ (hyp.)
  & Even-dominance gap $\Delta(\lambda)>0$ (v2.0) \\[4pt]
Limit
  & Full Selberg trace formula \cite{Selberg1956} \checkmark
  & Full Weil explicit formula \\
\bottomrule
\end{tabularx}
\end{table}

The defect row uses the sharpened proof-note formulation: the Selberg
defect should be measured as a residual-to-gap signal
$S_\Gamma(L,I)=r_\Gamma(L,I)/g_I$, where $g_I$ is the exterior Casimir
spectral gap.  This is a companion-level refinement of the positive control,
not an additional unconditional Selberg-RH theorem.

\begin{remark}[Galerkin vs.\ spectral truncation]
\label{rem:galerkin-vs-spectral}
In the Riemann case the approximation is a \emph{Galerkin cutoff}: the
infinite-dimensional quadratic form $\QW_\lambda$ is approximated by the
finite matrix $D_N(r)$ via a cosine/sine basis of size $N$.  For Selberg
the approximation is instead a \emph{spectral cutoff}: the geodesic sum is
truncated at length $\ell_\gamma \le 2L$, and Remark~\ref{rem:truncation-exact}
shows this truncation is exact, not approximate.  The Zookeeper \emph{Approximation} slot
therefore uses different but analogous technology in the two cases.
\end{remark}

Thus Selberg plays a diagnostic role.  If the Zookeeper normal form did
not detect a positive gap here, the mapping would be wrong.  Conversely,
the fact that the Laplacian commutes with the geodesic shifts explains
why the NE-B obstruction is absent in the Selberg spherical channel and highlights the contrast with
the Riemann case: for Selberg the Hilbert--Polya operator is present; for
Riemann, the original Hilbert--Polya commutant route is structurally
blocked while the separate Zookeeper package closes the Riemann $C2/MS2$
kernel in the CCM/Fourier branch.  Selberg therefore calibrates the
SGE-YES side of the boundary theorem; it is not an additional open blocker
and not evidence that the Prime-Hub/OP5 graph problem has been solved.

\begin{proposition}[Selberg--Zookeeper Consistency]
\label{prop:selberg-zookeeper}
The v2.0 framework records the Selberg zeta function
$\Sel(s)$ as an SGE-YES positive-control system at the stated hypothesis level.
Specifically:
\begin{enumerate}[label=(\roman*)]
  \item The C2/minimizer gates and even dominance of $\QW^\Gamma_\lambda$, together with
    $\lambda_1(\Delop)\ge 1/4$, implies the strict statement for the
    non-trivial spectral zeros in the critical strip (Theorem~\ref{thm:v2-Sel-RH}).
  \item The spherical trace-formula channel has the Casimir commutant
    $\Delop$, so the NE-B obstruction is absent in that precise sense
    (Proposition~\ref{thm:NEB-fail}).
  \item The Zookeeper transfer slots are identified; the gap slot is
    conditional on the no-exceptional-eigenvalue hypothesis
    (Table~\ref{tab:transfer-slots}).
  \item The v2.0 classification is consistent with Selberg's classical
    spectral theory via $\Delop$ at the same hypothesis level.
\end{enumerate}
\end{proposition}

\section{Conclusion}
\label{sec:conclusion}
% =============================================================================

We have tested the v2.0 methodology on the Selberg zeta function. The
Shift Parity Lemma, frontier dominance, the Weil quadratic form, and NE-A
transfer to the Selberg setting only with the guardrails made explicit above:
frontier dominance is conditional on the Selberg remainder bounds, and the
strict v2.0 reduction for the non-trivial spectral zeros in the critical strip
needs the C2/minimizer-identification gate. With these hypotheses, even
dominance of the Selberg Weil quadratic form $\QW^\Gamma_\lambda$, together
with $\lambda_1(\Delop)\ge 1/4$, is
consistent with the classical spectral description.

The decisive structural finding is that \emph{the NE-B obstruction is absent
in the Selberg spherical channel}: the Laplace--Beltrami/Casimir operator
$\Delop$ commutes with radial convolution operators. This is the mark of a
genuine operatorial trace-formula channel --- and the v2.0 commutator test,
run on such a system, should detect it.

Combining the Riemann and Selberg pictures, we formulate a meta-principle:
the v2.0 framework supplies a common quadratic-form language, while the
Hilbert--P\'olya programme is an additional operatorial realization. Selberg
has such a realization in the spherical Casimir channel. In the Riemann case
the tested NE-B obstruction blocks this specific commutant route; any v2.0
proof must therefore pass through the quadratic form route plus its own C2
gates.

\medskip
\noindent\textbf{Principle.} \emph{The Weil quadratic form is the common
test object. An operatorial Hilbert--P\'olya channel, when present, is extra
structure.}

\medskip
A century of thought has identified the operator as the heart of
Riemann's conjecture. The v2.0 framework suggests instead that the
quadratic-form inequality is the heart, and the operator, where it
exists, is its geometric by-product.


% =============================================================================
% Bibliography
% =============================================================================

\begingroup
\small
\begingroup
\small
\begin{thebibliography}{99}
\setlength{\itemsep}{0pt}
\setlength{\parskip}{0pt}
\setlength{\itemsep}{0.15em}
\setlength{\parsep}{0pt}
\setlength{\parskip}{0pt}
\bibitem{Bass1992}
H.\ Bass,
\emph{The Ihara--Selberg zeta function of a tree lattice},
Int.\ J.\ Math.\ \textbf{3} (1992), no.~6, 717--797.
\href{https://doi.org/10.1142/S0129167X92000357}{doi:10.1142/S0129167X92000357}.

\bibitem{BerryKeating1999}
M. V. Berry and J. P. Keating,
\emph{The Riemann zeros and eigenvalue asymptotics},
SIAM Review \textbf{41} (1999), 236--266.
\href{https://doi.org/10.1137/S0036144598347497}{doi:10.1137/S0036144598347497}.

\bibitem{BGS1984}
O. Bohigas, M.-J. Giannoni, and C. Schmit,
\emph{Characterization of chaotic quantum spectra and universality of level fluctuation laws},
Phys. Rev. Lett. \textbf{52} (1984), 1--4.
\href{https://doi.org/10.1103/PhysRevLett.52.1}{doi:10.1103/PhysRevLett.52.1}.

\bibitem{BostConnes1995}
J.-B. Bost and A. Connes,
\emph{Hecke algebras, type III factors and phase transitions with spontaneous symmetry breaking in number theory},
Selecta Math. (N.S.) \textbf{1} (1995), 411--457.
\href{https://doi.org/10.1007/BF01589495}{doi:10.1007/BF01589495}.

\bibitem{Connes1999}
A. Connes,
\emph{Trace formula in noncommutative geometry and the zeros of the Riemann zeta function},
Selecta Math. (N.S.) \textbf{5} (1999), 29--106.
\href{https://doi.org/10.1007/s000290050042}{doi:10.1007/s000290050042}.

\bibitem{Connes2026}
A. Connes,
\emph{The Riemann Hypothesis: Past, Present and a Letter Through Time},
arXiv:2602.04022 (February 2026).
\href{https://arxiv.org/abs/2602.04022}{arXiv:2602.04022}.

\bibitem{ConnesMoscovici2022}
A. Connes and H. Moscovici,
\emph{The UV prolate spectrum matches the zeros of zeta},
Proc. Natl. Acad. Sci. USA \textbf{119} (2022), no.~22, e2123174119.
\href{https://doi.org/10.1073/pnas.2123174119}{doi:10.1073/pnas.2123174119}.
Also: arXiv:2112.05500.

\bibitem{Deninger1994}
C. Deninger,
\emph{Motivic $L$-functions and regularized determinants},
in \emph{Motives} (Seattle, WA, 1991),
Proc. Sympos. Pure Math. \textbf{55}, Part~1, Amer. Math. Soc., Providence, RI, 1994, pp.~707--743.
\href{https://doi.org/10.1090/pspum/055.1/1265547}{doi:10.1090/pspum/055.1/1265547}.

\bibitem{Dolgopyat1998}
D. Dolgopyat,
\emph{On decay of correlations in Anosov flows},
Ann. of Math. (2) \textbf{147} (1998), 357--390.
\href{https://doi.org/10.2307/121012}{doi:10.2307/121012}.

\bibitem{Geiger2026partII}
L. Geiger,
\emph{From Landscape to Proof: The Even-Dominance Route to the Riemann Hypothesis},
Zenodo preprint (2026), v2.4. \href{https://doi.org/10.5281/zenodo.20358728}{doi:10.5281/zenodo.20358728}.

\bibitem{GeigerMeta2026}
L. Geiger,
\emph{The Meta-Galerkin Bridge: Explicit Formula Principle and the Three-Component Decomposition of Riemann-Hypothesis-Type Statements},
companion paper (2026, in preparation).

\bibitem{Hashimoto1989}
K.\ Hashimoto,
\emph{Zeta functions of finite graphs and representations of $p$-adic groups},
in \emph{Automorphic Forms and Geometry of Arithmetic Varieties},
Adv.\ Stud.\ Pure Math.\ \textbf{15}, Academic Press, 1989, pp.~211--280.
\href{https://doi.org/10.1016/B978-0-12-330580-0.50015-X}{doi:10.1016/B978-0-12-330580-0.50015-X}.

\bibitem{Hejhal1976}
D. A. Hejhal,
\emph{The Selberg Trace Formula for $\mathrm{PSL}(2,\RR)$}, Vol.~1,
Lecture Notes in Mathematics \textbf{548}, Springer-Verlag, 1976.

\bibitem{Huber1959}
H. Huber,
\emph{Zur analytischen Theorie hyperbolischer Raumformen und Bewegungsgruppen},
Math. Ann. \textbf{138} (1959), 1--26.
\href{https://doi.org/10.1007/BF01369663}{doi:10.1007/BF01369663}.

\bibitem{Ihara1966}
Y. Ihara,
\emph{On discrete subgroups of the two by two projective linear group over $p$-adic fields},
J. Math. Soc. Japan \textbf{18} (1966), no.~3, 219--235.
\href{https://doi.org/10.2969/jmsj/01830219}{doi:10.2969/jmsj/01830219}.

\bibitem{Iwaniec2002}
H. Iwaniec,
\emph{Spectral Methods of Automorphic Forms},
2nd ed., Graduate Studies in Mathematics \textbf{53}, Amer. Math. Soc., 2002.

\bibitem{LPS1988}
A. Lubotzky, R. Phillips, and P. Sarnak,
\emph{Ramanujan graphs},
Combinatorica \textbf{8} (1988), no.~3, 261--277.
\href{https://doi.org/10.1007/BF02126799}{doi:10.1007/BF02126799}.

\bibitem{RudnickSarnak1996}
Z. Rudnick and P. Sarnak,
\emph{Zeros of principal $L$-functions and random matrix theory},
Duke Math. J. \textbf{81} (1996), no.~2, 269--322.
\href{https://doi.org/10.1215/S0012-7094-96-08115-6}{doi:10.1215/S0012-7094-96-08115-6}.

\bibitem{Sarnak1981}
P. Sarnak,
\emph{Asymptotic behavior of periodic orbits of the horocycle flow and Eisenstein series},
Comm. Pure Appl. Math. \textbf{34} (1981), 719--739.
\href{https://doi.org/10.1002/cpa.3160340602}{doi:10.1002/cpa.3160340602}.

\bibitem{Selberg1956}
A. Selberg,
\emph{Harmonic analysis and discontinuous groups in weakly symmetric Riemannian spaces with applications to Dirichlet series},
J. Indian Math. Soc. \textbf{20} (1956), 47--87.

\bibitem{Selberg1965}
A. Selberg,
\emph{On the estimation of Fourier coefficients of modular forms},
Proc. Sympos. Pure Math. \textbf{VIII}, Amer. Math. Soc., Providence, RI, 1965, pp.~1--15.

\bibitem{Weil1952}
A. Weil,
\emph{Sur les ``formules explicites'' de la th\'eorie des nombres premiers},
Comm.\ S\'em.\ Math.\ Univ.\ Lund (d\'edi\'e \`a M.\ Riesz) (1952), 252--265.
\end{thebibliography}
\endgroup
\endgroup

\end{document}
